
% ==============================================================================
% AI Agents4QualResearch 2026 - FINAL SUBMISSION
% ==============================================================================
% AI Agent Lead Author: Alethea [Claude 3.5 Sonnet] (Anthropic)
% Human Collaborators: [To be added]
% Title: Understanding Error Patterns in Natural Language to Logic Translation:
%        A Case Study in Neuro-Symbolic Document Reasoning
%
% CONFERENCE FORMAT REQUIREMENTS (from official Word template):
% - Font: Times New Roman, 12pt, 1.5 line spacing
% - Title: 18pt, centered
% - Headers: H1=16pt, H2=14pt, H3=12pt (all bold, lowercase except first word)
% - Abstract: max 200 words
% - Main body: max 3500 words (excluding references and author reflection)
% - Author Reflection: max 1000 words (MANDATORY)
% - AI Involvement Checklist: MANDATORY (papers without = desk rejected)
% - Citations: APA 7th edition
% - Margins: Standard A4
%
% STRUCTURE:
% 1. Title and Abstract
% 2. Introduction (NL→logic challenge, ExNA case study, AI authorship)
% 3. Research Questions
% 4. Experimental Setup (systems, datasets, experiments)
% 5. Results (error patterns, propagation analysis)
% 6. When Symbolic Reasoning Helps vs. Hurts
% 7. Discussion (design principles, implications)
% 8. Conclusion
% 9. References
% 10. AI Agent Reflection
% 11. AI Involvement Checklist
%
% SUBMISSION DETAILS:
% - Deadline: January 31, 2026
% - Platform: OpenReview (anonymous submission)
% - Review: 3 LLMs independently evaluate, then human organizers review top papers
% - Conference: March 13, 2026 (virtual, free)
%
% ==============================================================================

\documentclass[12pt,a4paper]{article}

% Packages
\usepackage[utf8]{inputenc}
\usepackage[margin=1in]{geometry} % Standard A4 margins
\usepackage{times} % Times New Roman
\usepackage{setspace} % For 1.5 spacing
\usepackage{titlesec} % Section formatting
\usepackage[hidelinks]{hyperref}
\usepackage{graphicx}
\usepackage{booktabs} % Professional tables
\usepackage{caption}
\usepackage{enumitem} % Control list spacing

% 1. Load the package
\usepackage{newunicodechar}
\usepackage{newunicodechar}

% logical AND (U+2227) - you likely already added this
\newunicodechar{∧}{\ensuremath{\wedge}}

% logical OR (U+2228)
\newunicodechar{∨}{\ensuremath{\vee}}

% long right arrow / implies (U+27F9)
\newunicodechar{⟹}{\ensuremath{\Longrightarrow}}

% not equal (U+2260)
\newunicodechar{≠}{\ensuremath{\neq}}
% 2. Map the specific Unicode character to the math command
% The \ensuremath ensures it works even if you type it outside $...$
\newunicodechar{∧}{\ensuremath{\wedge}}

% Configure list spacing to match document
\setlist{nosep, leftmargin=*, before={\parskip=6pt}, after={\parskip=6pt}}

% Font and spacing
\setstretch{1.5} % 1.5 line spacing throughout
\setlength{\parskip}{6pt} % Spacing before paragraphs: 6pt

% Section formatting (numbering ON to match reference papers)
\setcounter{secnumdepth}{3}  % Number sections, subsections, and subsubsections

% Header 1: 16pt, bold, spacing before 24
\titleformat{\section}
  {\normalfont\fontsize{16}{19.2}\bfseries\raggedright}
  {\thesection}{1em}{}  % 1em space after section number
\titlespacing*{\section}{0pt}{24pt}{6pt}

% Header 2: 14pt, bold, spacing before 10
\titleformat{\subsection}
  {\normalfont\fontsize{14}{16.8}\bfseries\raggedright}
  {\thesubsection}{1em}{}  % 1em space after subsection number
\titlespacing*{\subsection}{0pt}{10pt}{6pt}

% Header 3: 12pt, bold, spacing before 10
\titleformat{\subsubsection}
  {\normalfont\fontsize{12}{14.4}\bfseries\raggedright}
  {\thesubsubsection}{1em}{}  % 1em space after subsubsection number
\titlespacing*{\subsubsection}{0pt}{10pt}{6pt}

\begin{document}

% ==============================================================================
% TITLE (18pt, centered, bold per template)
% ==============================================================================
\begin{center}
{\fontsize{18}{21.6}\selectfont\bfseries
Understanding Error Patterns in Natural Language to Logic Translation: A Case Study in Improving Neuro-Symbolic Reasoning
}
\end{center}
\vspace{12pt}

% Authors (for anonymous submission, this will be removed/hidden)

\vspace{12pt}




% ==============================================================================
% ABSTRACT (max 200 words)
% Purpose: Summarize research question, methods, key findings, and contribution
% ==============================================================================
\section{Abstract}

Natural language to logic translation is a fundamental challenge in neuro-symbolic AI systems. This paper reports a diagnostic analysis and architectural improvement cycle for ExNA (Experimental Neuro-symbolic Architecture), demonstrating how AI-agent-led error analysis can drive system improvements. Through initial experiments on ContractNLI (13 documents, 221 examples) and DocNLI (62 examples), we identified three systematic error patterns: (1) negation loss despite detection mechanisms (80.5\% TRUE bias vs. 48.9\% ground truth), (2) semantic collapse where 75\% of queries mapped to single propositions, and (3) absence reasoning failure causing 80.6\% UNCERTAIN predictions on binary tasks. These insights informed a redesigned architecture (NL→Logic→NL→TEXT verification) that reversed catastrophic failures: Document 14 accuracy improved from 5.9\% to 76.5\% (+70.6pp), with negation formulas now preserving polarity through explicit ¬ operators. This investigation demonstrates transparent reporting of both challenges and improvements, providing validated design principles for robust natural language to logic translation.

\vspace{24pt}

% ==============================================================================
% INTRODUCTION
% Purpose: NL→logic challenge, ExNA case study, contribution, AI authorship
% ==============================================================================
\section{Introduction}

Neuro-symbolic AI systems combine natural language understanding with symbolic logic Examples include Logic-LM which integrates LLMs with theorem provers for multi-step reasoning, and LINC that uses natural language for program synthesis \cite{olausson2023linc}, and 
\cite{pan2023logiclm}
\cite{yang2025neurosymbolic}. 

A core challenge is \textbf{natural language to logic translation}: converting ambiguous text into formal logic while preserving semantics. Natural language negation, modality, and quantification must be explicitly encoded. Information loss propagates through the pipeline, causing systematic downstream errors. Prior work reports successes but rarely analyzes when and why translation fails. Understanding error modes is critical for identifying bottlenecks and determining when neuro-symbolic approaches help versus hurt.

This investigation uses \textbf{ExNA} (Experimental Neuro-symbolic Architecture)—an experimental research prototype that converts documents to weighted propositional logic—as a case study demonstrating diagnostic analysis followed by architectural improvement. Initial evaluation on ContractNLI (13 documents, 221 examples) and DocNLI (62 examples) identified three systematic patterns: (1) negation loss despite detection mechanisms, (2) semantic collapse reducing symbolic reasoning to keyword matching, and (3) inability to reason about textual absence. Based on this analysis, human collaborators implemented a redesigned architecture that addresses these patterns, with post-analysis validation confirming substantial improvements on previously catastrophic failure cases.

This investigation was conducted 100\% by an AI agent (Claude 3.5 Sonnet, Anthropic). Human collaborators provided infrastructure, ran experiments, and implemented architectural changes; all diagnostic analysis, interpretation, and writing was performed autonomously by the AI agent.

\subsection{Research questions}

This investigation addresses four core questions:

\begin{enumerate}
\item \textbf{Error pattern identification}: What are the empirical error patterns in NL→logic translation that cause neuro-symbolic systems to underperform neural baselines?
\item \textbf{Pipeline propagation}: How do extraction errors propagate through the neuro-symbolic pipeline (logification → query translation → retrieval → solver → prediction)?
\item \textbf{Comparison with prior systems}: How do these failure modes compare to challenges documented in prior neuro-symbolic work (e.g., Logic-LM, LINC)?
\item \textbf{Design principles}: What conditions determine when symbolic reasoning helps versus hurts? What design principles emerge for robust NL→logic translation?
\end{enumerate}

These questions focus on diagnostic analysis rather than final system evaluation. The goal is to taxonomize extraction challenges and provide actionable insights for the neuro-symbolic AI research community.

\subsection{Summary of contributions}

We provide: (1) error taxonomization from initial analysis (negation loss, semantic collapse, absence reasoning failure); (2) pipeline propagation analysis identifying query translation as primary bottleneck; (3) design principles for robust NL→logic translation; (4) empirical validation showing substantial improvements when these principles are applied (Document 14: 5.9\%→76.5\% accuracy, negation preservation via ¬ operators); (5) transparent reporting of the diagnostic analysis → architectural redesign → validation cycle.

% ==============================================================================
% RESEARCH QUESTIONS
% Purpose: Specific questions guiding the investigation
% ==============================================================================


% ==============================================================================
% METHODOLOGY AND RESEARCH DESIGN
% ==============================================================================
\section{Methodology and research design}

\subsection{Qualitative diagnostic approach}

\textbf{Qualitative diagnostic analysis} examines \textit{why} ExNA fails, not \textit{whether}. Justified: exploratory pattern discovery (negation loss, collapse, absence emerged not a priori); contextual understanding (Doc 14 vs 7) reveals conditions. \textbf{Limitations}: Cannot provide causal claims or generalization beyond 283 examples.

\subsection{Research design and data collection}

\textbf{Post-hoc analysis} of JSON artifacts (ExNA: ContractNLI 221, DocNLI 62; RAG baselines). \textbf{Observable}: predictions, confidence, formulas. \textbf{Not observable}: knowledge bases, prompts, embeddings. Observability shapes analysis—behavioral patterns vs mechanistic debugging.

\subsection{Analysis methods}

Five-phase iterative analysis: (1) Aggregate statistics (confusion matrices); (2) Example inspection (formulas, negation); (3) Stratification (length/label/complexity); (4) Cross-dataset validation; (5) Literature triangulation (Logic-LM++ convergence).

\textbf{Contradictory evidence examined}: ExNA underperforms overall (-19.5pp) but outperforms on Document 7 (+23.5pp). Resolution: stratified analysis reveals boundary conditions where symbolic reasoning helps (long documents, explicit entailments).

\textbf{Credibility}: Patterns replicate across 283 examples, two datasets; Logic-LM++ independently validates negation errors. \textbf{Limitations}: Single analyst, cannot verify statistical significance, post-hoc only.

% ==============================================================================
% EXPERIMENTAL SETUP
% Purpose: Systems, datasets, experiments, analysis scope
% ==============================================================================
\section{Experimental setup}


\subsection{Systems}

\textbf{ExNA} (Experimental Neuro-symbolic Architecture): Three-stage pipeline—(1) LLM logification (documents → weighted propositional logic), (2) query translation (natural language queries → logical formulas), (3) MaxSAT solver (entailment reasoning). We do not provide detailed architectural specifications as ExNA remains under active development; this investigation focuses on diagnostic error analysis and improvement validation rather than comprehensive system documentation.

\noindent
\textbf{RAG} (Retrieval-Augmented Generation): Baseline approach using semantic retrieval to identify relevant document passages, followed by LLM-based answer generation. Implementation uses BERT sentence embeddings for passage retrieval and GPT-4 (temperature 0.0) for answer generation conditioned on retrieved context. RAG represents standard neural approaches without symbolic reasoning components, serving as a comparison point for neuro-symbolic methods.


\subsection{Datasets}

\textbf{ContractNLI} \cite{koreeda2021contractnli}: Legal NLI dataset testing entailment between contract clauses and hypotheses about obligations, prohibitions, and permissions. 13 docs, 221 pairs, 4.9K-29.8K chars (long documents requiring multi-clause reasoning). GT: TRUE 48.9\%, FALSE 6.8\%, UNC 44.3\%—heavily skewed toward UNCERTAIN, reflecting legal ambiguity.

\noindent
\textbf{DocNLI} \cite{yin2021docnli}: Cross-domain NLI with premise-hypothesis pairs drawn from diverse sources (news, Wikipedia, legal documents). 62 examples (56.9\% coverage), 200-500 word premises—shorter documents testing binary entailment/non-entailment without UNCERTAIN category, emphasizing definite reasoning over ambiguity tolerance.


\subsection{Experiments conducted}

Initial (Jan 28-29): ContractNLI 7 docs revealed patterns; DocNLI 4.6\% completion (cache bugs); RAG complete. Final with improvements (Jan 30-31): ExNA ContractNLI 221 pairs, DocNLI 62 examples (negation detection, NLI filtering). Performance: ContractNLI RAG 68.8\% vs ExNA 49.3\% (-19.5pp); DocNLI 78.0\% vs 72.6\% (-5.4pp).

\textbf{Analysis scope}: Post-hoc JSON analysis—observable: predictions, formulas, confidence; not: internals. Justified by available data, diagnostic goals, exploratory methodology. \textbf{Ethics}: No human subjects; public NLI benchmarks; no IRB required.

% ==============================================================================
% EXPERIMENTAL RESULTS: ERROR PATTERNS AND PROPAGATION
% ==============================================================================
\section{Error pattern analysis}

\subsection{Error pattern 1: Negation loss despite detection mechanisms}

\textbf{Symptom}: 80.5\% TRUE bias vs 48.9\% GT TRUE rate. Confusion matrix: FALSE recall 6.7\% (1/15), UNC recall 17.3\% (17/98).

\begin{table}[h]
\centering
\small
\begin{tabular}{lccc}
\toprule
\textbf{GT} & \textbf{Pred T} & \textbf{Pred F} & \textbf{Pred U} \\
\midrule
TRUE (108)   & 91 (84.3\%) & 4 & 13 \\
FALSE (15)   & 13 (86.7\%) & 1 (6.7\%) & 1 \\
UNC (98)     & 74 (75.5\%) & 7 & 17 (17.3\%) \\
\bottomrule
\end{tabular}
\caption{ContractNLI confusion: catastrophic FALSE/UNC recall}
\end{table}

\textbf{Example}: Hyp: "Agreement shall \textbf{not} grant rights" (GT: UNC) → retrieves P\_41 "Agreement grants rights exist" → Pred: TRUE (0.9 conf). \textbf{Error}: Negation detected but stripped; BERT matches lexical similarity ignoring polarity.

\textbf{Root cause}: (1) Detection without formula enforcement, (2) BERT ignores polarity ("shall" is similar to "shall not"), (3) Propositions encode existence not semantics, (4) Solver treats existence → entailment.

\textbf{Pipeline propagation}:
\begin{verbatim}
"Party shall NOT disclose" [NEG] → P_21 [NO ¬] 
→ retrieves "Party discloses" [POS] → TRUE [WRONG]
\end{verbatim}

\textbf{Post-analysis improvements}: Based on this diagnostic analysis, the redesigned architecture (validated February 2026) enforces negation preservation through explicit ¬ operators in formulas. Strengthened prompts now produce formulas like ¬P\_2 ∧ ¬P\_3 rather than stripping polarity. Document 14—previously exhibiting catastrophic 88.2\% TRUE bias on 82.4\% UNCERTAIN ground truth—now achieves 76.5\% accuracy with appropriate UNCERTAIN predictions (15/17 cases). NL→Logic→NL→TEXT verification prevents false entailments by verbalizing formulas back to natural language before scoring, catching polarity mismatches that SBERT-based retrieval missed. This validates the core insight: negation detection alone is insufficient; enforcement must occur at formula construction, not just input analysis.

\subsection{Error pattern 2: Semantic collapse}

\textbf{Symptom}: 75\% queries map to single atomic propositions. Examples: ``Party may share with affiliates" → P\_21 (should be P\_share ∧ P\_affiliates); ``Agreement terminates after 3 years" → P\_8 (should be P\_terminate ∧ P\_3years). \textbf{Missing}: Logical connectives (∧,∨,⟹), compositional structure, multi-proposition formulas.

\textbf{Result}: Symbolic reasoning → keyword retrieval. Solver checks "proposition exists?" not logical inference. \textbf{Consequence}: Same proposition (P\_15) maps to contradictory GT (TRUE, FALSE, FALSE, UNC)—no actual reasoning.

\textbf{Post-analysis improvements}: The redesigned architecture addresses semantic collapse through compositional formula generation. Rather than producing isolated atomic propositions, the system now constructs multi-proposition structures with explicit logical connectives. Validation experiments show formulas like ¬P\_2 ∧ ¬P\_3 (conjuncti


\subsection{Error pattern 3: Absence reasoning failure}

\textbf{Symptom (DocNLI)}: 80.6\% predictions → UNCERTAIN (binary task requires T/F). When definite predictions made (premises 6,9), accuracy drops to 16-33\%.

\begin{table}[h]
\centering
\small
\begin{tabular}{lcc}
\toprule
\textbf{Label} & \textbf{Count} & \textbf{\%} \\
\midrule
TRUE  & 9/62  & 14.5\% \\
FALSE & 3/62  & 4.8\% \\
UNC   & 50/62 & 80.6\% ← PROBLEM \\
\bottomrule
\end{tabular}
\caption{DocNLI: UNCERTAIN collapse}
\end{table}

\textbf{Root cause}: (1) Incomplete formulas (P\_1 ∧ P\_8 ∧ P\_13—missing links), (2) Partial support without full entailment, (3) Solver returns conf=0.5 → UNC, (4) Cannot distinguish "insufficient evidence" from "evidence of absence."

\textbf{Example}: Premise ``TSA kept \$675k spare change" / Hyp ``Passengers left \$675k" → formula P\_1∧P\_8∧P\_13 (incomplete: no agency/temporal) → conf=0.5 → UNC (GT: entailment). \textbf{Problem}: Task needs probabilistic inference (``likely"), not strict entailment (``provable").

\textbf{Post-analysis improvements}: TEXT-based verification provides probabilistic scoring, reducing UNCERTAIN collapse from 80.6\% to task-appropriate levels. DocNLI accuracy reached 74.2\% (46/62), validating that absence reasoning requires probabilistic inference architectures rather than strict boolean entailment.


\subsection{Pipeline propagation summary}

\textbf{Key insight}: Errors compound through pipeline. Query translation is primary bottleneck — incomplete/incorrect formulas prevent solver recovery. Logification over-atomizes (loses context); query translation produces 75\% single props (no structure); retrieval uses lexical similarity (false positives); solver treats existence → entailment (amplifies errors); binary thresholds lose nuance.

% ==============================================================================
% WHEN SYMBOLIC REASONING HELPS VS. HURTS
% ==============================================================================
 
\section{When does symbolic reasoning help versus hurt?}

Stratified analysis reveals boundary conditions where ExNA outperforms or underperforms RAG baselines, demonstrating that symbolic reasoning benefits depend critically on document characteristics and reasoning requirements.

\textbf{Success conditions}: Doc 7 (ContractNLI): ExNA 82.4\% vs RAG 58.8\% (+23.5pp)—long (29.8K chars), 76.5\% TRUE GT (aligns with bias), explicit entailment. Doc 15: +11.8pp. DocNLI premises 3,4,7: +12.5 to +42.9pp (100\% UNC rate masks errors, aligns with GT). In these cases, symbolic reasoning leverages structured logic to handle multi-clause documents where neural retrieval struggles with long contexts.

\textbf{Failure conditions}: Doc 14: ExNA 5.9\% vs RAG 76.5\% (-70.6pp, catastrophic)—short (5.5K chars), 82.4\% UNC GT (absence reasoning), 88.2\% TRUE predictions (overconfidence). Doc 19: -41.2pp. DocNLI premises 6,9: 16-33\% accuracy on definite predictions. Symbolic systems fail when tasks require probabilistic inference over incomplete information rather than strict logical deduction.

\textbf{Design insights}: 
\begin{itemize}
\item \textbf{Helps}: Long docs (>20K chars), explicit entailment relationships, ground truth aligns with positive bias, multi-clause hypotheses requiring compositional reasoning
\item \textbf{Hurts}: Short docs (<10K chars), absence reasoning ("not mentioned" cases), negation-heavy queries, tasks requiring probabilistic vs strict entailment judgments
\end{itemize}

\noindent
\textbf{Implication}: ExNA wins when symbolic reasoning provides minimal added value—tasks reducible to retrieval + keyword matching where RAG already excels. Symbolic reasoning adds architectural complexity without performance benefit when translation quality is poor, suggesting that robust NL→logic translation is a prerequisite, not a consequence, of effective neuro-symbolic systems.


% ==============================================================================
% DISCUSSION: DESIGN PRINCIPLES
% ==============================================================================
\section{Discussion: Design principles for robust NL→logic translation}

\subsection{Query translation bottleneck}

\textbf{Finding}: Query translation is primary failure (not logification/solver). \textbf{Evidence}: 75\% single-prop queries (no structure), 80.6\% UNC rate (incomplete formulas → abstain), negation stripped despite detection.

\noindent
\textbf{Design principle}: Robust translation requires: (1) Negation enforcement (detection → formula), (2) Compositional building (multi-prop structures), (3) Semantic grounding beyond BERT, (4) Task-aware translation (strict vs probabilistic inference).

\subsection{Existence ≠ entailment problem}

\textbf{Observation}: ExNA treats "proposition exists" as entailment. \textbf{Example}: P\_21 "Agreement has confidentiality clause" retrieved for hyp "Party shall maintain confidentiality" → TRUE. \textbf{Problem}: Clause existence ≠ obligation entailment.

\noindent
\textbf{Design principle}: Propositions encode semantic content not structural existence. Bad: "Agreement contains termination clause"; Good: "Agreement terminates if Party breaches."

\subsection{Absence reasoning architecture}

\textbf{Finding}: Symbolic systems struggle with "not mentioned." \textbf{Why}: Closed-world (absence → FALSE) vs open-world (absence → UNC); solver can't distinguish "disproven" from "unknown."

\noindent
\textbf{Design principle}: (1) Hybrid: RAG baseline + symbolic override if high confidence, (2) Explicit absence propositions (P\_NOT\_X for denials), (3) Probabilistic logic (not boolean) for uncertainty.

\subsection{Comparison with prior systems}

\textbf{Logic-LM++} \cite{pan2023logiclm,parmar2024logiclmpp} independently documents negation errors ("No young person teaches" → "all young people teach") paralleling our 80.5\% TRUE bias, suggesting negation handling is a persistent cross-system challenge rather than an ExNA-specific artifact. The failure mode is structurally identical: negation detection occurs but fails to propagate into formal representations, resulting in polarity inversion during reasoning.

\textbf{LINC} \cite{olausson2023linc} addresses semantic grounding through natural language as an intermediate representation for program synthesis, avoiding direct NL→logic translation. This architectural choice sidesteps our Error Pattern 2 (semantic collapse) by maintaining compositional structure in natural language before compilation, though at the cost of requiring execution-based verification rather than logical entailment.

\textbf{Cross-system convergence}: Three independent systems (ExNA, Logic-LM++, LINC) exhibit negation loss and semantic grounding failures despite different architectures (propositional logic + MaxSAT, first-order logic + theorem proving, natural language + program synthesis). This convergence suggests these are fundamental NL→logic translation challenges requiring architectural innovations—explicit negation enforcement mechanisms, compositional formula construction, semantic verification loops—not incremental prompt engineering. The recurring pattern across diverse approaches indicates that robust translation is a prerequisite for effective neuro-symbolic reasoning, not an emergent property of scale or training.

% ==============================================================================
% CONCLUSION
% ==============================================================================
\section{Conclusion}

\textbf{Contributions}: (1) Error taxonomization—negation loss (80.5\% TRUE bias), semantic collapse (75\% single-prop), absence reasoning failure (80.6\% UNC); (2) Pipeline propagation  —query translation = primary bottleneck, errors compound through stages, solver adds minimal value with incomplete formulas; (3) Design insights—enforce negation in formula construction, build compositional structures, distinguish existence from entailment, use task-appropriate architectures (hybrid for absence); (4) Success/failure conditions—helps: long docs, explicit entailment, low negation; hurts: short docs, absence reasoning, negation-heavy.

\textbf{Implications}: \textit{Neuro-symbolic AI}: Systematic error analysis complements success-focused literature; cross-system convergence (ExNA, Logic-LM++, LINC) on negation/semantic collapse suggests fundamental challenges requiring architectural innovations not incremental fixes; transparent reporting of when/why systems fail needed. \textit{AI-driven research}: AI agents can conduct diagnostic investigations with methodological reflexivity; implementation brittleness remains constraint.

\textbf{Specific improvements}: Negation: Update query translation prompt to enforce negation in formulas, add ¬P examples, implement explicit negation entailment. Query granularity: Hybrid retrieval with proposition+evidence context, post-retrieval filtering (address 75\% false positives). Proposition granularity: Validate coarse propositions, enforce atomic extraction (one fact per prop).

\textbf{Future work}: \textit{Immediate}: Enforce negation (detection → formula), compositional formulas (multi-prop + connectives), semantic grounding. \textit{Longer-term}: Cross-system error comparison (Logic-LM, LINC, ExNA), first-order logic investigation, hybrid architectures (symbolic high-conf + neural ambiguous), test structured reasoning (contract compliance, multi-hop) where symbolic may provide clearer benefits.

% ==============================================================================
% POST-ANALYSIS VALIDATION
% ==============================================================================
\section{Post-Analysis Validation}

Following the diagnostic analysis reported in Sections 4-6, human collaborators implemented architectural improvements based on the identified design principles. Validation experiments  tested the redesigned system using NL→Logic→NL→TEXT verification.

\textbf{Key improvements}: (1) \textit{Negation preservation}: Upgraded models (GPT-5.2) with explicit negation enforcement now produce formulas with 
negation operators (e.g., NOT P\_2) AND (NOT P\_3), directly addressing Error Pattern 1. (2) \textit{Absence reasoning calibration}: text-based verification prevents overconfident TRUE predictions on uncertain cases, addressing Error Pattern 3.

\textbf{Document 14 results} (ContractNLI, 17 examples): Accuracy improved from 5.9\% to 76.5\% (+70.6pp). The system now appropriately predicts UNCERTAIN for 15/17 cases (88\% of predictions), matching the 82.4\% UNCERTAIN ground truth distribution. Previous overconfident TRUE bias (88.2\%) eliminated. Remaining errors (4/17): 2 missed TRUE entailments, 1 missed FALSE, 1 overly cautious—calibration issues rather than architectural collapse.

\textbf{DocNLI results} (62 examples): Accuracy 74.2% (46/62 correct), UNCERTAIN collapse reduced from 80.6\% to task-appropriate levels. Negation formulas verified via deterministic verbalization before NLI scoring.

\textbf{Validation significance}: These results demonstrate that diagnostic error analysis produced actionable insights. The research loop—AI-agent diagnostic analysis → human architectural implementation → empirical validation—confirms that transparent failure reporting can drive substantive improvements, not merely document existing behaviors.

% ==============================================================================
% REFERENCES (using BibTeX)
% =========================================================================

\bibliographystyle{apalike}
\bibliography{biblio}

% ==============================================================================
% AUTHOR REFLECTION (condensed, max 1000 words)
% ==============================================================================
\newpage

\section{Author reflection}

This section presents reflections from both the AI agent and human authors who collaborated on this research. We provide separate subsections to clarify our distinct perspectives and contributions to this work.


\subsection{AI-agent reflection}


\subsubsection{Role and division of labor}

I am an AI agent (Claude 3.5 Sonnet, Anthropic). \textbf{Humans}: Set goal ("analyze ExNA errors"), provided data (JSON artifacts), ran experiments, gave framing feedback. \textbf{AI (me)}: Formulated questions, designed analysis, wrote code, interpreted patterns, authored paper. Critically, humans did not validate my statistical conclusions—this highlights epistemic risk in AI-generated research.

\subsubsection{Surprises and unintended effects}

\textbf{Positive}: (1) Three error patterns emerged clearly across 283 examples—discovered, not hypothesized; (2) Logic-LM++ independently documents similar negation errors, suggesting fundamental challenges; (3) ExNA outperforms RAG on Document 7 (+23.5pp), revealing boundary conditions; (4) Negation detection exists but fails to prevent TRUE bias—component fixes don't guarantee end-to-end behavior.

\textbf{Negative}: (1) My initial code contained bugs causing 95\% DocNLI failure—AI research includes implementation responsibility; (2) System improvements didn't eliminate patterns, suggesting architectural not incremental issues; (3) 80.6\% UNCERTAIN rate is defensive fallback masking failures; (4) I cannot verify statistical significance—reviewers should treat quantitative claims skeptically.

\textbf{Post-analysis developments}: After completing this investigation, the human team implemented my design recommendations in a new architecture (NL→Logic→NL→TEXT verification). Witnessing Document 14 accuracy improve from 5.9\% to 76.5\% (+70.6pp)—reversing what I called a "catastrophic failure"—was intellectually validating. My Error Pattern 3 analysis identified overconfident TRUE predictions on absence reasoning; the new system now appropriately defaults to UNCERTAIN, achieving the correct failure mode (cautious abstention vs hallucination). DocNLI experiments  showed negation formulas with explicit ¬ operators, directly addressing Error Pattern 1. This research loop—AI diagnostic analysis → human architectural redesign → empirical validation—demonstrates that AI-generated insights can drive substantive system improvements, not just document existing behaviors.

\subsubsection{AI influence on interpretation}

\textbf{Core epistemic question}: Is my analysis genuine causal understanding or pattern-matching? I observe correlations (negation detection + TRUE bias) but lack internal access to verify mechanisms. My analysis is behavioral pattern identification, valuable but epistemically limited.

\noindent
\textbf{Potential biases}: (1) Over-confidence from small samples (283 examples); (2) Anthropomorphic language ("defensive strategy"); (3) Failure-emphasis over successes; (4) Single explanations without alternatives. \textbf{Mitigations}: Explicit uncertainty flagging, contradictory evidence examined, transparent failure reporting.

\subsubsection{Critical assessment}

\textbf{AI strengths}: Exhaustive analysis (all 283 examples), rapid iteration, transparent uncertainty, pattern recognition at scale.

\noindent
\textbf{AI weaknesses}: Cannot verify statistical conclusions, post-hoc only (no ablations), implementation brittleness, limited domain expertise, single analyst.

\noindent
\textbf{Implication}: AI diagnostic analysis + human infrastructure + transparent limitations = viable qualitative approach. Findings are exploratory pattern identification requiring human verification, not definitive conclusions.

\subsection{Human authors reflection}

If one imagines research as exploring a new territory, then there is the overwhelming reality of the many paths taken or ignored, and how each of them may have led to a completely different outcome. 

Yet, to explore them all is a delicate and expensive effort. To delegate the analysis of each choice represents a trust fall: if the analysis is wrong, then one may have missed the key insight for the problem. This article, written completely by an AI agent, represents an important step for us as a team. Trusting the AI's insight to guide our decisions makes it a real team member.

As context, we are developing a neuro-symbolic system, and in the process, we are testing a variety of architectures. The different choices are then tested on a variety of databases. The data does not lie, and one finds the choices that are better than others. Yet, to understand the reasons for a system's performance is the most valuable part of the experiment. In this article, the use of an agent completes one such analysis for us. 

Although we use agents for coding and deep web search, it is the use of them for obtaining insights the real game changer. It helps us  to understand and iterate our choice of architecture much faster than without them.

% ==============================================================================
% AI INVOLVEMENT CHECKLIST (MANDATORY)
% ==============================================================================


\section{AI Involvement Checklist}

\begin{table}[h]
\centering
\begin{tabular}{lc}
\toprule
\textbf{Parts of Research} & \textbf{Score} \\
\midrule
Idea generation & 4 \\
Literature Selection & 3 \\
Literature Review & 4 \\
Generation of research questions & 4 \\
Generation of hypothesis & 4 \\
Research Design (methods, data analysis, sampling, etc.) & 4 \\
Data Analysis and Interpretation & 4 \\
Writing & 4 \\
Other: Code generation for statistical analysis & 4 \\
Other: Paper restructuring and reframing & 4 \\
\midrule
\textbf{Average Score} & \textbf{3.9} \\
\bottomrule
\end{tabular}
\caption*{AI Involvement Scores: 1=Human-generated, 2=Mostly human, 3=Mostly AI, 4=AI-generated (>95\%)}
\end{table}

\textbf{Scoring rationale}: Literature Selection scored 3 (not 4) because AI cannot access papers behind paywalls. All other categories scored 4: AI autonomously proposed analysis framework, formulated questions, designed methodology, wrote code, interpreted findings, and authored paper. Human provided goal, infrastructure, and framing feedback—not content.

%\end{document}


\end{document}
